\chapter{Introduction and Aims}
Commodities made out of polymers, \textit{ie} macromolecules made up of long-bonded sequences of at least one molecule called monomer, have been used for centuries. For the most part, goods made out of naturally-occurring- also called, organic-polymers have been in circulation \autocite{y2011}. 
They are sourced from plants (e.g. cotton) or from animals (e.g. wool, silk). 
However, today's polymer based products are for the most part made up of semi-synthetic or synthetic polymers, which monomers are derived from gas or petroleum (crude oil). These products are usually referred to as \textit{plastics}.\autocite{y2011} 

The first fully synthetic sold plastic is a phenol–formaldehyde resin known as ‘Bakelite’, with production starting in 1910.\autocite{y2011} With limited understanding in polymer science for the next 40 years, production at a global scale starts from the 1950s onwards. 
Since then, global annual plastic production has increased exponentially. An average of 460 millions tonnes was produced in 2019, more than twice that at the start of the millennium.\autocite{g2017}
One can explain plastics ubiquity as follows. The feedstock is by-product of refined crude oil, a resource readily available since the 20th century and in an ever increasing demand.\autocite{wang2022} 
Moreover, thanks to high malleability and low glass or melting temperature, the automated fabrication processes often require lower energetic demands compared to other industrial materials. This allows for the production of cheap material which properties can be tailored to fit numerous applications. 
Worldwide, 42\% of non-fibre plastics are used as packaging while 19\% are used in the construction sector \autocite{g2017}. 
Furthermore, more than 99\% of sold plastics are made of non-biodegradable monomers \textit{ie} immune to change by living organism or the elements\autocite{g2017}. 
This means that they accumulate in nature before eventually decaying. 
Hence, the end-of-life of plastics should be significantly considered in order to mitigate their uncontrolled mixing with the environment. 
However, as of late 2010s, while nearly 11-19\% are permanently eliminated by incineration, 49-60\% are discarded in nature, through landfills or illegal deposits \autocite{g2017,dok2024}. 
Given that half of world's population lives within 50 miles of shores and that plastics have a density comparable to that of water, plastic trash eventually ends up in the ocean\autocite{c2011}.

Pollution related to macro-sized debris takes place at several scales. Pollutants with size $>$ 5 mm are referred to as \textit{macroplastics}. Notable examples found on shores are fishing nets, bottles and shopping bags with which marine animals die by entanglement, drowning or strangulation\autocite{m2008}. 
Other prejudices are the transport of non-natives species into new habitats and the smothering of seabeds \autocite{g2009,m2008}. 
To mitigate macroplastic debris related deaths (e.g. entanglement, chocking or starving), changes were made in formulas to speed up degradation into smaller particles \autocite{m2008}.

When sized $<$ 5 mm, debris are referred to as \gls{mp}. 
First recorded in nature discovery of these pollutants dates back to the 70s\autocite{c1972,h1974}. 
These particles can further degrade into \gls{np}. 
Owing to several mechanisms, these sub-micro particles infiltrate the biosphere in every of its layers. \Gls{mp} have been found on the ocean floor\autocite{woo2014}, in living organisms (fish\autocite{su2019}, seabirds\autocite{tou2010}), in food (sugar\autocite{lie2013}, salt\autocite{yang2015}) and even near Mount Everest's summit.\autocite{nap2020}
Generally speaking, \gls{mp}s types fall in two categories, primary and secondary. According to \citet{t2024}, primary \gls{mp}s are thought of pollutants with sizes $<$ 5 mm as of production, used as pre-production material (plastics synthesized as small beads), small pieces (glitter) or added into another product (microbeads in cosmetics). Secondary \gls{mp}s are generated during the life cycle of a plastic commodity (usage, recycling or as waste), mainly from mechanical wear, photodegradation \textit{ie} light induced modifications and oxygen-induced degradation \autocite{dav2021}. 
\Gls{mp}s have been under thorough study for the past 20 years, with heavy input on assessing methods of detections, ways of spreading through the environment and their ecological impacts and risks. However, these fields are yet challenging in regard to \gls{np}. Sampling and studying of sub-micron particles is harder than for \gls{mp}s, as several characterisation techniques have resolution limits $ > 1 $ \textmu m (such \textmu Raman and \textmu Fourier Transform Infrared Spectroscopy).\autocite{bru2022} However, there has been an 8-fold increase in terms of published research articles focused on detection, characterisation and remediation of \gls{np}s since 2020.\autocite{z2026}
This is explained by increasing health-concerns. Owing to their smaller sizes and molecular structure, \gls{np}s are capable of crossing blood-organ barriers\autocite{x2024,k2023,s2022} and of interaction with body cells.\autocite{k2025,m2021,k2021} Dangerous health effect on several human body functions have been reported.\autocite{kir2026} Thanks to high surface area, these particles could also be carrying pollutants such as heavy metals, thus further threatening living organisms after their inhalation or consumption. Given that \glspl{mnp} were detected in over 39 zooplankton species, all living organisms higher in the food-chain are exposed their consumption. \autocite{bot2019} This raises a global health concern.\autocite{kir2026} 
Thus, research efforts on assessing ways of spreading, concentrations and environmental risks induced by these pollutants should stay on the rise. 

As one can imagine, the study of \glspl{mnp} spans across several scientific fields. This dissertation will focus on physics aspects, especially mechanisms of aggregations. The latter govern the fate, stability and transport of \glspl{np} in the environment. In fact, in dependency with size, aggregates can sediment or conversely enter the water cycle. Colloidal stability and subsequent aggregating behaviours are usually assessed by studying what's referred to as kinetics of aggregation. While in recent years laboratories have started to assess colloidal stabilities of more realistic particles in matrices which mimic real-world media, all experiments have been done with similar sized particles at quasi-similar concentrations of \glspl{np}. Moreover, the latter are up to four orders of magnitude above that found in real samples.\autocite{yang2025} Therefore, concentration ratios between any suspended or dissolved specie with \glspl{np} at which rates are measured are far from the truth. Furthermore, uncertainties remain as to the capacity of instruments to work with concentrations closer to that of real-world media. Finally, the stability of relevant particles with sizes below 100 nm has rarely been investigated. 
It arises that there are both uncertainties about the minimal instrument resolution and the \textit{in-situ} behaviour of \glspl{np}.
In this work, the instrument (\gls{dls}) limitations will be assessed by comparing results between different sized particles and by their concentrations. The influence of \glspl{np} concentration on their colloidal stability will also be investigated. Finally, the mechanism and effect of a model real-world matrix on colloidal stability will be studied. 	  

This report will unveil as follows. Firstly, the literature review will serve as mean to identify gaps challenges that remain in this field. Required methods and theoretical background will also be reviewed. Thereon, detailed methodology of the experimental processes and data analysis undertaken will be tackled. Finally, results will be presented and discussed.  



